\newpage
\section{Significance of the Study}
\label{ch:significance}

This study may contribute to the development of more reliable automated methods for assessing rice pollen viability, particularly when pollen grains overlap or occur in clusters. Previous research has demonstrated the potential of image analysis and deep learning for pollen identification and viability assessment, but performance can vary with pollen type and image conditions~\cite{go2019,zhang2023,tan2022,li2026pvdt}. By evaluating a YOLOv8n-seg model on stained rice pollen images, this study will provide evidence on the usefulness and limitations of instance segmentation for this specific task.

For rice researchers and plant breeders, an effective automated method could help reduce the amount of manual counting needed when examining large numbers of pollen grains. By producing grain-level classifications and estimating the proportion of viable pollen, the proposed approach may offer a useful supporting tool for pollen assessment. It is intended to complement, rather than replace, established laboratory procedures; its practical value will depend on the results of the evaluation and the quality of the available images.

For researchers in computer vision and agricultural applications of artificial intelligence, this study will examine how a pretrained instance-segmentation model performs on a specialized microscopy task with overlapping objects. The comparison with the image-analysis protocol of Go et al.~\cite{go2019}, together with separate evaluation of images with and without lumped pollen, will help clarify whether the proposed approach addresses the limitations identified in the existing method. The use of cross-validation and comparison with manual counts will also provide a structured basis for assessing detection, classification, and viability-estimation performance.

Finally, the findings may serve as a reference for future studies on automated pollen analysis. If the proposed method proves effective, subsequent work could investigate additional rice genotypes, larger and more varied image datasets, other segmentation models, or validation against complementary biological measures. If its performance remains limited, the results may still identify challenges that need to be addressed before deep-learning-based pollen viability assessment can be applied more broadly.
